\chapter{Bounded-Suboptimal Search: ECBS}
\label{ch:ch10}
% Function names for pseudocode are prefixed with "Ecbs" to stay unique
% across the book.
\SetKwFunction{EcbsFocalSearch}{FocalSearch}
\SetKwFunction{EcbsLowLevel}{FocalSpaceTimeAStar}
\SetKwFunction{EcbsHighLevel}{ECBS}

\begin{objectives}
\begin{itemize}
  \item Explain why optimal \cbs stops scaling after a dozen interacting drones and what a suboptimality factor $w$ buys: a plan whose cost is provably at most $w$ times the optimum, found in a small fraction of the time.
  \item Run focal search by hand: keep \Open ordered by $\fcost$, keep \Focal as the band of open nodes with $\fcost \le w\,\fcost_{\min}$, choose inside the band by a second heuristic, and prove that the result is $w$-suboptimal.
  \item Trace \ecbs on a small instance: a low level that returns a path \emph{and} a lower bound, a high level that runs focal search over the constraint tree, and the lower-bound bookkeeping that ties the two together.
  \item State and prove the theorem $\cost \le w\,C^*$ for \ecbs, and explain why the algorithm terminates on solvable instances.
  \item Implement \ecbs in \python, benchmark it against \cbs, choose $w$ from the measured speed--quality trade-off, and recognise the three classic implementation mistakes.
  \item Describe what \emph{EECBS} adds: explicit estimation search and online-learned estimates.
\end{itemize}
\end{objectives}

% ---------------------------------------------------------------------
\section{Why this matters for a drone swarm}
\label{sec:ch10-motivation}

Picture an inspection swarm of sixteen drones on a warehouse roof cut into an
$8 \times 8$ grid of cells, a tenth of them blocked by ventilation units. The
operator presses \emph{start} and expects the drones to move within a second.
Conflict-based search (\cbs) from \cref{ch:ch09} produces the cheapest
conflict-free plan, but on random instances of exactly this kind, in the
experiment of \cref{sec:ch10-benchmark}, it finds that plan for only two of ten
instances within ten seconds; the other eight are still expanding
constraint-tree nodes when the limit strikes. The reason is not a weak
implementation. Every conflict that \cbs resolves splits the constraint tree in
two, and the number of nodes with a cost below the optimum grows exponentially
with the number of conflicts that the optimal plan has to untangle. Optimal
search pays for a certificate that the operator never asked for: the proof that
no cheaper plan exists.

The operator asked for something weaker and far cheaper: a plan that is
\emph{good enough}, with a number attached that says how good. This is what a
\textbf{bounded-suboptimal}\index{bounded suboptimality} algorithm delivers.
You choose a factor $w \ge 1$, and the algorithm guarantees a solution whose
sum of costs is at most $w$ times the optimal sum $C^*$, without ever computing
$C^*$. Enhanced \cbs (\ecbs)\index{ECBS} by Barer, Sharon, Stern and
Felner~\cite{barer2014ecbs} is the bounded-suboptimal version of \cbs. On the
same sixteen-drone instances, \ecbs with $w = 1.5$ solves all ten in about
$30$ milliseconds each, and the plans it returns are on average $7\,\%$ more
expensive than the optimum, far below the guaranteed $50\,\%$. The mechanism
behind it, \textbf{focal search}\index{focal search}, is thirty years older
than \cbs: Pearl and Kim's $\astar_\eps$~\cite{pearl1982studies} keeps the
ordinary \astar queue as a ruler and lets a second heuristic pick, among the
nodes that the ruler declares to be within the bound, the one that seems to
lead to a solution fastest. In \ecbs that second heuristic is simply the number
of conflicts.

For the hybrid architecture of \cref{ch:ch24} this chapter settles a design
decision. The global planner that produces the nominal, time-indexed paths of
the swarm is \cbs or \ecbs, and for more than a handful of drones it is \ecbs:
the local avoidance layer perturbs the nominal paths anyway, the replanning
layer calls the global planner again whenever a drone has been pushed off its
route, and a planner that sometimes needs minutes is useless in that loop. The
bound $w$ is the knob that trades flight time for planning time, and the lower
bound that \ecbs computes on the way tells you, after each run, how much of the
optimum you gave away. Week~5 of the study plan (\cref{ch:appA}) asks you to
add \ecbs to your \cbs solver and to measure exactly this trade-off; its
milestone is that you can \emph{measure the cost of faster suboptimal
solutions}.

\paragraph{Roadmap.} \Cref{sec:ch10-intuition} explains focal search in
pictures. \Cref{sec:ch10-problem} fixes the notation: bounded suboptimality,
\Focal, per-agent lower bounds. \Cref{sec:ch10-algorithm} gives three
algorithms: generic focal search, the \ecbs low level in space-time and the
\ecbs high level, each with a walkthrough. \Cref{sec:ch10-example} runs \cbs
and \ecbs on the same three-agent instance and compares cost, tree size and
bound for several values of $w$. \Cref{sec:ch10-properties} proves the
$w$-suboptimality of focal search and of \ecbs and discusses completeness.
\Cref{sec:ch10-variants} covers the choice of $w$, the family of suboptimal
\cbs variants, a benchmark against \cbs and the ideas of EECBS.
\Cref{sec:ch10-implementation,sec:ch10-drone} give the implementation notes,
the pitfalls and the place of \ecbs in the drone system.

% ---------------------------------------------------------------------
\section{The idea in plain words}
\label{sec:ch10-intuition}

\subsection{A band instead of a point}

\astar (\cref{ch:ch04}) always expands the open node with the smallest
$\fcost = \gcost + \hcost$. That rule is what makes it optimal, and it is also
what makes it stubborn: even when a node with a slightly larger $\fcost$ would
finish the search in three steps, \astar first works through every node whose
$\fcost$ is smaller. Weighted \astar (\cref{ch:ch04}) loosens the rule by
inflating the heuristic, $\fcost_w = \gcost + w\,\hcost$. It is fast, and it
is $w$-suboptimal, but its only notion of ``promising'' is the inflated
$\hcost$: it cannot use any other information about a node.

Focal search separates the two jobs. \Open stays exactly as in \astar, ordered
by an admissible $\fcost$. Its smallest value, $\fcost_{\min}$, is a lower
bound on the optimal cost $C^*$, because some node on an optimal path is
always open and has $\fcost \le C^*$ (\cref{sec:ch10-properties}). Therefore
\emph{every} open node with
\begin{equation}
  \fcost(n) \le w\,\fcost_{\min}
  \label{eq:ch10-band}
\end{equation}
can be expanded without breaking the bound: if it is a goal, its cost is at
most $w\,\fcost_{\min} \le w\,C^*$. These nodes form the
\textbf{focal list}\index{FOCAL list} \Focal, a band of width $w$ at the front
of \Open (\cref{fig:ch10-idea}). Inside the band any rule may pick the node to
expand. The rule is a second heuristic $\hcost_{\Focal}$ that need not be
admissible and need not even estimate cost: it may estimate the number of
steps to a solution, or, in \ecbs, the number of conflicts. As long as the
choice stays inside the band, the solution is within the factor $w$ of the
optimum, whatever $\hcost_{\Focal}$ says.

\begin{figure}[tb]
  \centering
  \inputfigure{ch10/idea}
  \caption{Focal search. The open nodes lie on the $\fcost$ axis; the band
  between $\fcost_{\min}$ and $w\,\fcost_{\min}$ (here $w = 1.3$) is \Focal.
  \astar would expand the leftmost node. Focal search expands the node inside
  the band with the smallest secondary value $\hcost_{\Focal}$, here the node
  with $\fcost = 12$ and $\hcost_{\Focal} = 0$. Nodes outside the band are
  ignored for now even if their $\hcost_{\Focal}$ is zero; they enter the band
  when $\fcost_{\min}$ grows.}
  \label{fig:ch10-idea}
\end{figure}

\subsection{Two levels, one bound}

\cbs has two searches, and both can be made focal. The \emph{low level} plans
one drone in space-time under a set of constraints. Its \Open is ordered by
$\fcost = t + \hcost(v)$ with the exact grid distance as $\hcost$, and its
\Focal prefers partial paths that collide with fewer of the other drones'
current paths. The result is a path that is at most $w$ times longer than the
drone's own constrained optimum and that tends to create fewer conflicts for
the high level to resolve. The \emph{high level} searches the constraint tree.
Its \Open is ordered by a lower bound on the cost of the best solution below
each node, and its \Focal prefers nodes whose paths contain fewer conflicts,
because a node with no conflicts \emph{is} a solution.

The subtle part is the lower bound. In \cbs the cost of a node is a lower
bound on the cost of every solution below it, because each path is optimal
for its constraints. In \ecbs the paths are not optimal, so the cost of a node
proves nothing. The fix is the piece of bookkeeping that gives the algorithm
its correctness: the low level returns, next to the path, the value
$\fcost_{\min}$ it saw when it stopped, which is a lower bound on the
constrained optimum of that drone. A node stores one such bound per drone,
their sum $\mathrm{LB}(N)$ orders \Open, and the smallest $\mathrm{LB}(N)$ in
\Open, called $\mathrm{LB}$, is a lower bound on $C^*$. \Focal is the set of
open nodes with $\cost(N) \le w\,\mathrm{LB}$. When a conflict-free node is
taken from \Focal, its cost is at most $w\,\mathrm{LB} \le w\,C^*$.

\begin{keyidea}
Keep the admissible ordering of \Open as a ruler and expand any open node whose
$\fcost$ is within the factor $w$ of the smallest one, choosing among them by
a second heuristic that need not be admissible. \ecbs does this at both
levels with the number of conflicts as the second heuristic, and it makes the
high-level ruler honest by letting every low-level search return a lower
bound together with its path.
\end{keyidea}

% ---------------------------------------------------------------------
\section{Problem statement and notation}
\label{sec:ch10-problem}

We use the MAPF notation of \cref{ch:ch07}. An instance consists of a
4-connected grid with blocked cells, $k$ agents $a_1, \dots, a_k$ with distinct
start cells $s_i$ and distinct goal cells $\gamma_i$, unit-time move and wait
actions, and the stay-at-target convention. A path $\pi_i$ lists the cell of
agent $a_i$ at every time step; $\pi_i[t] = \gamma_i$ for all $t$ after the
final arrival, and the cost $\cost(\pi_i)$ is the time of that final arrival.
A solution $\Pi = (\pi_1, \dots, \pi_k)$ must be free of vertex conflicts
$\langle i, j, v, t\rangle$ and swapping conflicts $\langle i, j, u, v, t\rangle$;
its cost is the sum of costs $\cost(\Pi) = \sum_i \cost(\pi_i)$, and $C^*$ is
the smallest cost of any solution. Cells are written $(r, c)$ with row $r$
counted from the top and column $c$ from the left. A vertex constraint
$\langle i, v, t\rangle$ forbids agent $a_i$ to be at $v$ at time $t$; an edge
constraint $\langle i, u, v, t\rangle$ forbids it to move from $u$ to $v$
between $t$ and $t+1$. A constraint-tree (CT) node $N$ holds a set of
constraints $N.\mathcal{C}_i$ for every agent, one path $N.\pi_i$ per agent
that respects $N.\mathcal{C}_i$, and the resulting $\cost(N) = \sum_i
\cost(N.\pi_i)$. This is the setting of \cbs in \cref{ch:ch09}.

\begin{definition}[Bounded-suboptimal algorithm]
\label{def:ch10-bounded}
Let $w \ge 1$. A solution $\Pi$ is \textbf{$w$-suboptimal}\index{w-suboptimal
solution@$w$-suboptimal solution} if $\cost(\Pi) \le w\,C^*$. An algorithm
with parameter $w$ is \textbf{bounded-suboptimal}\index{bounded
suboptimality!definition} if, on every solvable instance, it returns a
$w$-suboptimal solution. The number $w$ is the \textbf{suboptimality
factor}\index{suboptimality factor}. For $w = 1$ the algorithm is optimal.
\end{definition}

The definition asks for a guarantee, not for equality: the returned solution
may be much better than $w\,C^*$, and usually is
(\cref{sec:ch10-benchmark}). Weighted \astar with an admissible heuristic is a
bounded-suboptimal algorithm for single-agent search (\cref{ch:ch04}). So is
focal search, which we now define for a general search problem: a finite
graph with non-negative edge costs, a start node $s$, a goal test, and an
admissible heuristic $\hcost$, exactly as in \cref{ch:ch04}.

\begin{definition}[Focal search]
\label{def:ch10-focal}
Let $\hcost$ be admissible, let $\fcost(n) = \gcost(n) + \hcost(n)$, and let
$\hcost_{\Focal} \colon V \to \R$ be any function, the \textbf{secondary
heuristic}\index{secondary heuristic}\index{heuristic!secondary (focal)}. At
every step of a best-first search with open list \Open, write
$\fcost_{\min} = \min_{n \in \Open} \fcost(n)$ and
\begin{equation}
  \Focal = \set{n \in \Open : \fcost(n) \le w\,\fcost_{\min}}.
  \label{eq:ch10-focal}
\end{equation}
\textbf{Focal search}\index{focal search!definition} with factor $w$ expands, at
every step, a node of \Focal with the smallest $\hcost_{\Focal}$. Pearl and
Kim~\cite{pearl1982studies} call it $\astar_\eps$ with $w = 1 + \eps$.
\index{A*-epsilon@$\astar_\eps$}
\end{definition}

Note the two heuristics have different duties. $\hcost$ must be admissible,
because $\fcost_{\min}$ has to be a true lower bound; $\hcost_{\Focal}$ may be
anything, because it only chooses inside a band that is safe anyway. Setting
$w = 1$ turns \Focal into the set of nodes with $\fcost = \fcost_{\min}$ and
focal search into \astar with $\hcost_{\Focal}$ as tie-breaker. For \ecbs the
secondary heuristic is a count of conflicts, at both levels.

\begin{definition}[Conflict heuristic]
\label{def:ch10-hc}
For a CT node $N$, $h_c(N)$\index{conflict heuristic $h_c$} is the number of
vertex and swapping conflicts among the paths of $N$. For a low-level state
$(v, t)$ of agent $a_i$, reached by a partial path from $(s_i, 0)$, $h_c(v, t)$
is the number of conflicts that this partial path has with the current paths
of the other agents up to time $t$; if $(v, t)$ is a goal state, the conflicts
that $a_i$ would have while parked at $\gamma_i$ from time $t$ on are added.
\end{definition}

\begin{definition}[Lower bounds of a CT node]
\label{def:ch10-lb}
Let $C_i^*(N)$ be the cost of a cheapest path of agent $a_i$ that respects
$N.\mathcal{C}_i$ (ignoring the other agents). A \textbf{per-agent lower
bound}\index{lower bound!per agent} is a number $\mathrm{LB}_i(N) \le
C_i^*(N)$ stored in $N$. The \textbf{lower bound of the node}\index{lower
bound!of a CT node} is $\mathrm{LB}(N) = \sum_{i=1}^{k} \mathrm{LB}_i(N)$,
and the \textbf{global lower bound}\index{lower bound!global LB} of the high
level is
\begin{equation}
  \mathrm{LB} = \min_{N \in \Open} \mathrm{LB}(N).
  \label{eq:ch10-LB}
\end{equation}
\end{definition}

$\mathrm{LB}(N)$ bounds the cost of every solution in the subtree below $N$,
because every such solution respects the constraints of $N$ and the paths of
different agents are bounded separately. \Cref{sec:ch10-properties} shows that
$\mathrm{LB} \le C^*$ at all times, which is what the high-level \Focal
\begin{equation}
  \Focal = \set{N \in \Open : \cost(N) \le w\,\mathrm{LB}}
  \label{eq:ch10-focal-high}
\end{equation}
relies on.

% ---------------------------------------------------------------------
\section{The algorithm}
\label{sec:ch10-algorithm}

\subsection{Focal search in general}
\label{sec:ch10-focal-alg}

\Cref{alg:ch10-focal} is focal search on an arbitrary graph, written as a
graph search with re-opening exactly like \astar in \cref{ch:ch04}; only the
choice rule differs.

\begin{algorithm}[htb]
\caption{Focal search ($\astar_\eps$) on a graph}
\label{alg:ch10-focal}
\KwIn{graph with costs $c \ge 0$, start $s$, goal test, admissible $\hcost$, secondary heuristic $\hcost_{\Focal}$, factor $w \ge 1$}
\KwOut{a path from $s$ to a goal of cost at most $w\,C^*$, or \emph{failure}}
\Fn{\EcbsFocalSearch{$s$, $\hcost$, $\hcost_{\Focal}$, $w$}}{
  $\gcost(s) \gets 0$; parent$(s) \gets$ \KwNil; \Open $\gets \set{s}$; \Closed $\gets \emptyset$\;
  \While{\Open $\ne \emptyset$}{
    $\fcost_{\min} \gets \min_{n \in \Open} \fcost(n)$\label{alg:ch10-focal:fmin}\;
    \Focal $\gets \set{n \in \Open : \fcost(n) \le w\,\fcost_{\min}}$\label{alg:ch10-focal:band}\;
    $n \gets \argmin_{n' \in \Focal} \hcost_{\Focal}(n')$\tcp*{ties: smaller $\fcost$, then larger $\gcost$}\label{alg:ch10-focal:pop}
    \If{$n$ is a goal}{\KwRet path to $n$ along the parent pointers\label{alg:ch10-focal:goal}\;}
    move $n$ from \Open to \Closed\;
    \ForEach{$(n', c(n, n')) \in \Succ(n)$}{
      \If{$\gcost(n) + c(n, n') < \gcost(n')$}{\label{alg:ch10-focal:relax}
        $\gcost(n') \gets \gcost(n) + c(n, n')$; parent$(n') \gets n$\;
        put $n'$ into \Open, removing it from \Closed if it was there\tcp*{re-opening}\label{alg:ch10-focal:reopen}
      }
    }
  }
  \KwRet \emph{failure}\;
}
\end{algorithm}

\paragraph{Walkthrough.}
Line~\ref{alg:ch10-focal:fmin} reads the ruler: the smallest $\fcost$ in
\Open. Line~\ref{alg:ch10-focal:band} marks the band; in an implementation
\Focal is not recomputed but maintained incrementally
(\cref{sec:ch10-implementation}). Line~\ref{alg:ch10-focal:pop} is the only
place where $\hcost_{\Focal}$ is consulted. Ties are broken towards the
smaller $\fcost$ and then the larger $\gcost$, so that with $w = 1$ the
algorithm degenerates to \astar with the tie-breaking of \cref{ch:ch04}. The
goal test at line~\ref{alg:ch10-focal:goal} happens when the node is chosen,
never when it is generated, for the same reason as in \astar: the proof of the
bound in \cref{sec:ch10-properties} argues about the moment of selection. The
relaxation at line~\ref{alg:ch10-focal:relax} is unchanged, and re-opening
(line~\ref{alg:ch10-focal:reopen}) is needed in general: because focal search
expands nodes out of $\fcost$ order, a node may be expanded with a $\gcost$
that is not yet optimal even when $\hcost$ is consistent, and a cheaper path to
it may turn up later. Note that $\fcost_{\min}$ never decreases when $\hcost$
is consistent: expanding a node generates only successors with a larger or
equal $\fcost$, and the node with the smallest $\fcost$ always lies inside the
band, so once it is gone the minimum can only rise.

\subsection{The low level: focal search in space-time}
\label{sec:ch10-lowlevel}

The low level of \ecbs is the space-time \astar of \cref{ch:ch04} with the
choice rule of \cref{alg:ch10-focal}. A state is a pair $(v, t)$; the cost of
reaching it is $g = t$; the heuristic $\hcost(v)$ is the exact grid distance
from $v$ to the goal, computed once by a backward breadth-first search and
cached for all time layers and all constrained re-runs. This heuristic is
consistent and ignores the other agents, so it never overestimates. Successors
are the wait action $(v, t+1)$ and the moves $(v', t+1)$ to free neighbours,
minus those forbidden by a vertex or edge constraint. Two details of
\cref{ch:ch04} carry over. The goal test accepts $(\gamma_i, t)$ only if $t$ is
later than every vertex constraint on $\gamma_i$, written $t > t_\gamma$, so
that the agent can stay at its goal for ever. And the search is cut at a time
horizon $T_{\max} = H + 1 + D$, where $H$ is the latest constrained time and
$D$ the largest grid distance to the goal, which keeps it finite and does not
exclude any optimal constrained path.

\begin{algorithm}[htb]
\caption{The \ecbs low level: focal space-time \astar with a lower bound}
\label{alg:ch10-lowlevel}
\KwIn{agent $a_i$ with start $s_i$ and goal $\gamma_i$, constraint set $\mathcal{C}_i$, current paths $\Pi_{-i}$ of the other agents, factor $w \ge 1$}
\KwOut{a path $\pi_i$ respecting $\mathcal{C}_i$ and a lower bound $\mathrm{lb}_i$ with $\mathrm{lb}_i \le C_i^*(\mathcal{C}_i)$ and $\cost(\pi_i) \le w\,\mathrm{lb}_i$; or \emph{failure}}
\Fn{\EcbsLowLevel{$a_i$, $\mathcal{C}_i$, $\Pi_{-i}$, $w$}}{
  $\hcost(v) \gets$ grid distance from $v$ to $\gamma_i$ for all free cells $v$ (backward BFS, cached)\label{alg:ch10-lowlevel:h}\;
  $t_\gamma \gets$ largest $t$ with $\langle i, \gamma_i, t\rangle \in \mathcal{C}_i$, or $-1$; $T_{\max} \gets H + 1 + D$\;
  $n_0 \gets (s_i, 0)$; $h_c(n_0) \gets$ conflicts of $n_0$ with $\Pi_{-i}$; \Open $\gets \set{n_0}$; \Closed $\gets \emptyset$\;
  \While{\Open $\ne \emptyset$}{
    $\fcost_{\min} \gets \min_{n \in \Open} \fcost(n)$; \Focal $\gets \set{n \in \Open : \fcost(n) \le w\,\fcost_{\min}}$\label{alg:ch10-lowlevel:band}\;
    $(v, t) \gets \argmin_{n \in \Focal} h_c(n)$ (ties: smaller $\fcost$, then larger $t$); move it to \Closed\label{alg:ch10-lowlevel:pop}\;
    \If{$v = \gamma_i$ \KwAnd $t > t_\gamma$}{\KwRet (path to $(v, t)$, $\fcost_{\min}$)\label{alg:ch10-lowlevel:goal}\;}
    \If{$t + 1 > T_{\max}$}{\KwContinue}
    \ForEach{$v' \in \set{v} \cup \Succ(v)$ with $\langle i, v', t+1\rangle \notin \mathcal{C}_i$ and $\langle i, v, v', t\rangle \notin \mathcal{C}_i$}{
      \If{$(v', t+1)$ was generated before}{\KwContinue\tcp*{same cost $t+1$: first path wins}\label{alg:ch10-lowlevel:dup}}
      $h_c(v', t+1) \gets h_c(v, t) + \abs{\set{j : \pi_j[t+1] = v'}} + \abs{\set{j : \pi_j[t] = v',\ \pi_j[t+1] = v}}$\label{alg:ch10-lowlevel:hc}\;
      \If{$v' = \gamma_i$ \KwAnd $t + 1 > t_\gamma$}{$h_c(v', t+1) \gets h_c(v', t+1) + \abs{\set{(j, t') : t' > t+1,\ \pi_j[t'] = \gamma_i}}$\label{alg:ch10-lowlevel:park}\;}
      $\fcost(v', t+1) \gets t + 1 + \hcost(v')$; parent$(v', t+1) \gets (v, t)$; put $(v', t+1)$ into \Open\label{alg:ch10-lowlevel:push}\;
    }
  }
  \KwRet \emph{failure}\;
}
\end{algorithm}

\paragraph{Walkthrough.}
Line~\ref{alg:ch10-lowlevel:h} prepares the heuristic; in \cref{ch:ch09} you
saw that the same table serves every low-level call for the same agent. The
band at line~\ref{alg:ch10-lowlevel:band} and the choice at
line~\ref{alg:ch10-lowlevel:pop} are \cref{alg:ch10-focal} with
$\hcost_{\Focal} = h_c$. Line~\ref{alg:ch10-lowlevel:goal} is the addition
that makes \ecbs work: the search returns the current $\fcost_{\min}$ together
with the path. At that moment the goal state itself is still counted in
\Open, so $\fcost_{\min} \le \fcost(\gamma_i, t) = t$; \cref{sec:ch10-properties}
shows that $\fcost_{\min}$ is also a lower bound on the optimal constrained
cost, and that $t \le w\,\fcost_{\min}$ because the goal was taken from the
band. Line~\ref{alg:ch10-lowlevel:dup} is the duplicate test. In space-time
with unit costs, every path to $(v', t+1)$ costs exactly $t+1$, so a second
visit can never be cheaper and re-opening is never needed; the first path to
a state wins, even if a later path to the same state would have fewer
conflicts. This simplification is standard, and \cref{sec:ch10-example} shows
one of its consequences. Line~\ref{alg:ch10-lowlevel:hc} counts the conflicts
that the step adds: the other agents that occupy $v'$ at time $t+1$ (a vertex
conflict) and those that move from $v'$ to $v$ while we move from $v$ to $v'$
(a swapping conflict); under stay-at-target, $\pi_j[t]$ is the goal of $a_j$
for every $t$ after its arrival. Line~\ref{alg:ch10-lowlevel:park} makes the
count of a goal state exact: an agent that parks at its goal from time $t+1$
on collides with every later visit of another agent to that cell, and the high
level must know about these conflicts too. Line~\ref{alg:ch10-lowlevel:push}
pushes the successor into \Open; whether it also belongs to \Focal follows from
\cref{eq:ch10-focal} and is decided by the implementation as described in
\cref{sec:ch10-implementation}.

\begin{figure}[tb]
  \centering
  \inputfigure{ch10/focal-lowlevel}
  \caption{The low level on a $3 \times 4$ grid. Agent B walks down column~1
  and parks at $(2,1)$; agent A must get from $(1,0)$ to $(1,3)$. The
  shortest path (dashed) costs~3 but meets B at $(1,1)$ at $t = 1$. With
  $w = 1$ the band contains only states with $\fcost = 3$ and the dashed path
  is returned with one conflict. With $w = 1.5$ the band admits
  $\fcost \le 4.5$; the state ``wait at $(1,0)$ until $t = 1$'' has
  $\fcost = 4$ and no conflict, so it is expanded first and the solid path of
  cost~4 is returned. Both runs return the lower bound $\fcost_{\min} = 3$.}
  \label{fig:ch10-focal-lowlevel}
\end{figure}

\begin{example}[The low level on a $3 \times 4$ grid]
\label{ex:ch10-lowlevel}
\Cref{fig:ch10-focal-lowlevel} shows a $3 \times 4$ grid without obstacles.
Agent B's path is fixed: $(0,1) \to (1,1) \to (2,1)$, where it parks. Agent A
starts at $(1,0)$ with goal $(1,3)$, so $\hcost(1,0) = 3$ and
$\fcost_{\min} = 3$ at the start. \Cref{tab:ch10-lowlevel-trace} traces
\cref{alg:ch10-lowlevel} for A with $w = 1.5$; the instance is
\code{low\_level\_example()} in \code{ch10\_ecbs.py} and the numbers are
checked by its self-test. In step~1 the start is expanded. Its successor
$(1,1)$ at $t = 1$ has $\fcost = 3$ but $h_c = 1$, because B is there at
$t = 1$; the wait $(1,0)$ at $t = 1$ has $\fcost = 4 \le 1.5 \cdot 3$ and
$h_c = 0$; the moves to $(0,0)$ and $(2,0)$ have $\fcost = 5 > 4.5$ and stay
outside the band. Step~2 expands the wait, because its $h_c$ is smaller,
although its $\fcost$ is larger. From there the search moves along row~1 one
step behind B; every state has $\fcost = 4$ and $h_c = 0$, and in step~5 the
goal state $(1,3)$ at $t = 4$ is taken from the band. The returned path costs
$4$, has no conflict, and comes with the lower bound $\fcost_{\min} = 3$,
which is the cost of the dashed path. Run with $w = 1$, the same code expands
$(1,1)$ at $t = 1$ in step~2 and returns the dashed path with one conflict
after three expansions. The threshold between the two behaviours is
$w = 4/3$ (\cref{exr:ch10-lowlevel-threshold}).
\end{example}

\begin{table}[tb]
  \centering\small
  \caption{Trace of \cref{alg:ch10-lowlevel} on \cref{ex:ch10-lowlevel} with
  $w = 1.5$. States are written cell$@t$. The last two columns show
  \Focal after the step, with $h_c/\fcost$ for each state, and the remaining
  open states with their $\fcost$; $\fcost_{\min}$ stays $3$ throughout,
  so the band is $\fcost \le 4.5$.}
  \label{tab:ch10-lowlevel-trace}
  \begin{tabular}{@{}rlrrp{4.0cm}p{4.6cm}@{}}
  \toprule
  Step & Expanded & $h_c$ & $\fcost$ & \Focal after the step ($h_c/\fcost$) & Open, outside the band ($\fcost$)\\
  \midrule
  1 & $(1,0)@0$ & 0 & 3 & $(1,1)@1{:}\,1/3$, $(1,0)@1{:}\,0/4$ & $(0,0)@1{:}5$, $(2,0)@1{:}5$\\
  2 & $(1,0)@1$ & 0 & 4 & $(1,1)@1{:}\,1/3$, $(1,1)@2{:}\,0/4$ & $(0,0)@1{:}5$, $(2,0)@1{:}5$, $(1,0)@2{:}5$, $(0,0)@2{:}6$, $(2,0)@2{:}6$\\
  3 & $(1,1)@2$ & 0 & 4 & $(1,1)@1{:}\,1/3$, $(1,2)@3{:}\,0/4$ & as before plus $(1,1)@3{:}5$, $(0,1)@3{:}6$, $(2,1)@3{:}6$, $(1,0)@3{:}6$\\
  4 & $(1,2)@3$ & 0 & 4 & $(1,1)@1{:}\,1/3$, $(1,3)@4{:}\,0/4$ & as before plus $(1,2)@4{:}5$, $(0,2)@4{:}6$, $(2,2)@4{:}6$, $(1,1)@4{:}6$\\
  5 & $(1,3)@4$ & 0 & 4 & goal: return path of cost 4, $\mathrm{lb} = 3$ & \\
  \bottomrule
  \end{tabular}
\end{table}

\subsection{The high level: focal search over the constraint tree}
\label{sec:ch10-highlevel}

\Cref{alg:ch10-ecbs} is the high level. Compared with \cbs, three things
change: every node carries per-agent lower bounds, \Open is ordered by
$\mathrm{LB}(N)$ instead of $\cost(N)$, and the node to expand is chosen from
\Focal by $h_c$.

\begin{algorithm}[htb]
\caption{\ecbs: the high level}
\label{alg:ch10-ecbs}
\KwIn{MAPF instance with agents $a_1, \dots, a_k$, factor $w \ge 1$}
\KwOut{a conflict-free solution $\Pi$ with $\cost(\Pi) \le w\,C^*$, or \emph{failure}}
\Fn{\EcbsHighLevel{instance, $w$}}{
  $R.\mathcal{C}_i \gets \emptyset$ for $i = 1, \dots, k$\;
  \ForEach{$i = 1, \dots, k$}{$(R.\pi_i, R.\mathrm{lb}_i) \gets$ \EcbsLowLevel{$a_i$, $\emptyset$, $(R.\pi_1, \dots, R.\pi_{i-1})$, $w$}\label{alg:ch10-ecbs:root}\;}
  $\cost(R) \gets \sum_i \cost(R.\pi_i)$; $\mathrm{LB}(R) \gets \sum_i R.\mathrm{lb}_i$; $h_c(R) \gets$ number of conflicts among $R$'s paths\;
  \Open $\gets \set{R}$\;
  \While{\Open $\ne \emptyset$}{
    $\mathrm{LB} \gets \min_{N \in \Open} \mathrm{LB}(N)$; \Focal $\gets \set{N \in \Open : \cost(N) \le w\,\mathrm{LB}}$\label{alg:ch10-ecbs:band}\;
    $N \gets \argmin_{N' \in \Focal} h_c(N')$ (ties: smaller cost, then older node); remove $N$ from \Open\label{alg:ch10-ecbs:pop}\;
    \If{$h_c(N) = 0$}{\KwRet $N$'s paths\tcp*{$\cost(N) \le w\,\mathrm{LB} \le w\,C^*$}\label{alg:ch10-ecbs:goal}}
    $\chi \gets$ the earliest conflict among $N$'s paths, between agents $a_i$ and $a_j$\label{alg:ch10-ecbs:conflict}\;
    \ForEach{$a \in \set{a_i, a_j}$, with $\chi_a$ the constraint on $a$ that resolves $\chi$}{
      $N' \gets$ copy of $N$; $N'.\mathcal{C}_a \gets N.\mathcal{C}_a \cup \set{\chi_a}$\label{alg:ch10-ecbs:child}\;
      $(\pi, \mathrm{lb}) \gets$ \EcbsLowLevel{$a$, $N'.\mathcal{C}_a$, paths of the other agents in $N'$, $w$}\label{alg:ch10-ecbs:replan}\;
      \If{the low level failed}{\KwContinue\tcp*{no path under these constraints: prune}}
      $N'.\pi_a \gets \pi$; $N'.\mathrm{lb}_a \gets \max(N.\mathrm{lb}_a, \mathrm{lb})$\label{alg:ch10-ecbs:lb}\;
      recompute $\cost(N')$, $\mathrm{LB}(N')$ and $h_c(N')$; insert $N'$ into \Open\label{alg:ch10-ecbs:insert}\;
    }
  }
  \KwRet \emph{failure}\;
}
\end{algorithm}

\paragraph{Walkthrough.}
Line~\ref{alg:ch10-ecbs:root} builds the root. We plan the agents one after
another and hand each low-level call the paths planned so far, so that cheap
conflict avoidance happens already at the root: agent $a_i$ steers around
$a_1, \dots, a_{i-1}$ where it can within its band, and ignores
$a_{i+1}, \dots, a_k$, whose paths do not exist yet. This is prioritized
planning (\cref{ch:ch08}) used as an initialisation, but nothing depends on
it: the root of plain \cbs, with independent shortest paths, would do as well.
Line~\ref{alg:ch10-ecbs:band} computes the global lower bound
\cref{eq:ch10-LB} and the band \cref{eq:ch10-focal-high};
line~\ref{alg:ch10-ecbs:pop} picks the node with the fewest conflicts inside
it. The goal test at line~\ref{alg:ch10-ecbs:goal} is the same as in \cbs, and
the comment states the guarantee proved in \cref{sec:ch10-properties}.
Lines~\ref{alg:ch10-ecbs:conflict} to~\ref{alg:ch10-ecbs:insert} are the
\cbs split: the earliest conflict is resolved by one constraint on each of the
two agents, and each child re-plans only the constrained agent. The other
agents keep their paths \emph{and their lower bounds}, which remain valid
because their constraint sets did not change.

Line~\ref{alg:ch10-ecbs:lb} is the bookkeeping. The replanned agent's bound is
the larger of two valid lower bounds: the bound it had in the parent, which
stays valid because adding a constraint cannot make the constrained optimum
cheaper, and the $\fcost_{\min}$ that the new low-level search returned. Barer
et al.\ use the returned value alone; taking the maximum is a safe tightening
(\cref{sec:ch10-properties}) and often lifts the global $\mathrm{LB}$ a little
earlier. Either way, $\mathrm{LB}(N')$ can be larger than $\cost$ of the
parent and, unlike the cost, it never decreases along a branch. The cost of a
child, on the other hand, may go down: the replanned path is only within
$w$ of the agent's optimum, and the parent's path may have been a poor one.
This is why \Open must be ordered by $\mathrm{LB}(N)$ and not by cost.

\paragraph{The two rulers.}
It helps to keep the two bands apart. The low level's band is
$\fcost \le w\,\fcost_{\min}$ inside \emph{one} agent's space-time search; it
guarantees $\cost(\pi_a) \le w\,\mathrm{lb}_a$ for every path that ever enters
a CT node. The high level's band is $\cost(N) \le w\,\mathrm{LB}$ over
\emph{all} agents and over the whole tree. The low-level guarantee is what
keeps the high-level band non-empty: the node $N_{\min}$ that attains
$\mathrm{LB}$ has $\cost(N_{\min}) = \sum_a \cost(\pi_a) \le w \sum_a
\mathrm{lb}_a = w\,\mathrm{LB}(N_{\min}) = w\,\mathrm{LB}$, so it always
belongs to \Focal. The high-level guarantee, $\cost \le w\,\mathrm{LB} \le
w\,C^*$, is the one the user sees.
